\par\Needspace{6\baselineskip}\medskip
\hypertarget{native-125-article092}{}\pdfbookmark[2]{92.}{article-092}
\begin{center}92.\end{center}
Obwohl die Gleichungen 15 und $15^{\ast}$ in manchen Fällen drei reelle Wurzeln haben können, ist stets eindeutig, welche davon für unsere Aufgabe zu wählen ist. Da $h$ offenbar positiv ist, folgt aus der Theorie der Gleichungen leicht, dass Gleichung 15 genau eine positive Wurzel und daneben entweder zwei imaginäre oder zwei negative Wurzeln hat. Da $y=\dfrac{m}{\BareRoot{(l+x)}}$ notwendig positiv ist, bleibt hier offenbar keine Unsicherheit \SourceLeaf{126}bestehen. Was Gleichung $15^{\ast}$ betrifft, so bemerken wir zunächst, dass $L$ notwendig größer als 1 ist. Das lässt sich leicht zeigen, indem man die in Art. 89 angegebene Gleichung in die Form $L=1+\dfrac{\cos\tfrac12 f^2}{-\cos f}+\dfrac{\operatorname{tang}2\omega^2}{-\cos f}$ bringt. Setzt man ferner in Gleichung $12^{\ast}$ für $M$ den Ausdruck $Y\BareRoot{(L-x)}$ ein, so erhält man $Y+1=(L-x)X$ und damit
\[
\begin{aligned}
Y+1>(1-x)X>{}&\tfrac43+\dfrac{4}{3\mathbin{.}5}x+\dfrac{4\mathbin{.}6}{3\mathbin{.}5\mathbin{.}7}xx
\\[4pt]
&+\dfrac{4\mathbin{.}6\mathbin{.}8}{3\mathbin{.}5\mathbin{.}7\mathbin{.}9}x^3+\text{usw.}>\tfrac43,
\end{aligned}
\]
also $Y>\tfrac13$. Setzen wir also $Y=\tfrac13+Y'$, so muss $Y'$ positiv sein, und Gleichung $15^{\ast}$ geht über in
\[
Y'^3+2Y'Y'+(1-H)Y'+\tfrac{4}{27}-\tfrac29H=0,
\]
Diese Gleichung kann, wie sich aus der Theorie der Gleichungen leicht zeigen lässt, höchstens eine positive Wurzel haben. Daraus folgt, dass Gleichung $15^{\ast}$ genau eine Wurzel hat, die größer als $\tfrac13$ ist\footnote{Sofern wir voraussetzen, dass die Aufgabe tatsächlich lösbar ist.}. Diese Wurzel ist für unsere Aufgabe zu wählen, die übrigen bleiben außer Betracht\HistoricalNoteDouble{Handschriftliche Bemerkung:] Die Gleichung $15^{\ast}$ hat
\par 1) Eine reelle negative Wurzel, wenn $H$ zwischen den Grenzen 0 und $\dfrac{+1+\BareRoot{5}}{6}$ liegt, nebst zwei imaginären.
\par 2) Drei reelle Wurzeln, darunter eine positive, wenn $H$ zwischen $\dfrac{+1-\BareRoot{5}}{6}$ und 0 liegt.
\par 3) Eine reelle positive Wurzel und zwei imaginäre für $H$ zwischen $-\infty$ und $\dfrac{+1-\BareRoot{5}}{6}\mathpunct{\cdot}$
\par 4) Drei reelle Wurzeln, darunter eine negative für $H$ zwischen $\dfrac{+1+\BareRoot{5}}{6}$ und $+\infty$.
\par Offenbar kann also nur von Fall 4 hier die Rede sein, wo sich zwei positive Wurzeln finden. Eine davon ist hier stets kleiner als $\dfrac{-1+\BareRoot{5}}{6}\mathpunct{\raisebox{.45ex}{,}}$ die andere größer. Daher kommt nur die letztere infrage.
\par Nach dem im Text angegebenen Verfahren ist klar, dass die letzte Gleichung nur dann zwei positive Wurzeln haben könnte, wenn zugleich $1-H$ negativ und $\tfrac{4}{27}-\tfrac29H$ positiv wäre. Das ist offenbar unmöglich, da $\tfrac{4}{27}-\tfrac29H=\tfrac29(1-H)-\tfrac{2}{27}$.
}.
